% lem-outer-crossing: limit-shapes.tex:578-583 (label lem:outer-crossing)
\paragraph{Paper excerpt.}
\begin{verbatim}
\begin{lemma}\label{lem:outer-crossing}
Let $\alpha>0$. There exists $C(d,\drift,\gauge,\alpha,p)<\infty$ with the following property. Let $q\in\R^d$ satisfy $|q|\leq\Lambda_\gauge$, and consider a path of the walk that satisfies~\eqref{eq:vector}, and~\eqref{eq:linear-mart} for this~$q$ and every $a\in\Lambda_\gauge\Z\cap[0,\Lambda_\gauge n]$. Let $x_*$ be one of the sites $X_0,\ldots,X_n$ with $q\cdot X_j\leq T\coloneqq q\cdot x_*$ for $0\leq j\leq n$, and let $h>0$ be such that $q\cdot\xi(X_j)\geq\alpha$ for every $0\leq j\leq n$ with $q\cdot X_j>T-h$. For $b\geq0$ let $L_b\coloneqq1+\max\{\ell_n(x):x\in\Z^d,\ q\cdot x\geq b\}$. Then every $b\geq0$ satisfies
\begin{equation}\label{eq:outer-crossing}
T\leq\max\{T-h,b\}+CL_b\log n\, .
\end{equation}
\end{lemma}
\end{verbatim}
\paragraph{Lean statement.}
\begin{verbatim}
theorem CERW.Frozen.outer_crossing {d : ℕ} (hd : 2 ≤ d) :
    ∀ Ψ : EuclideanSpace ℝ (Fin d) → ℝ, CERW.IsNorm Ψ →
    ∀ ε : ℝ, 0 < ε → (∀ i : Fin d, ε * Ψ (CERW.coordVec i) < 1 / (d : ℝ)) →
    ∀ p : ℝ, 0 < p → ∃ C₁ : ℝ, 0 < C₁ ∧
      (∀ ξ : Site d → EuclideanSpace ℝ (Fin d),
        (∀ x : Site d, x ≠ 0 → CERW.IsSubgradient Ψ (CERW.toSpace x) (ξ x)) → ξ 0 = 0 →
      ∀ {Ω : Type u} [MeasurableSpace Ω] (μ : Measure Ω) [IsProbabilityMeasure μ]
        (X : ℕ → Ω → Site d), CERW.IsDriftCERW μ ε ξ X → ∀ n : ℕ, 2 ≤ n →
        let Λ : ℝ := CERW.normMax Ψ
        let Z : Ω → ℕ → EuclideanSpace ℝ (Fin d) := fun ω t => CERW.toSpace (X t ω) + ε •
          ∑ j ∈ Finset.range t,
            if X j ω ∉ CERW.departureRange (X · ω) j then ξ (X j ω) else 0
        let E : Set Ω := {ω | (∀ s t : ℕ, s < t → t ≤ n →
            ‖Z ω t - Z ω s‖ ≤ C₁ * Real.sqrt ((t - s : ℝ) * Real.log n)) ∧
          ∀ q ∈ ((LatticeProb.ballFinset d (n : ℝ)).filter (fun x => x ≠ 0)).image
              (fun x => (Λ / ‖CERW.toSpace x‖) • CERW.toSpace x),
          ∀ k : ℕ, k ≤ n →
            |CERW.dynkinMart (CERW.driftStepProb d ε ξ)
                (fun z => max (inner ℝ q (CERW.toSpace z) - k * Λ) 0) (X · ω) n|
              ≤ C₁ * (Real.sqrt (Real.log n * ∑ z ∈ (CERW.departureRange (X · ω) n).filter
                    (fun z => k * Λ - Λ < inner ℝ q (CERW.toSpace z)),
                    (CERW.localTime (X · ω) n z : ℝ)) + Real.log n)}
        μ Eᶜ ≤ ENNReal.ofReal (C₁ * (n : ℝ) ^ (-p))) ∧
      ∀ α : ℝ, 0 < α → ∃ C : ℝ, 0 < C ∧
      ∀ ξ : Site d → EuclideanSpace ℝ (Fin d),
        (∀ x : Site d, x ≠ 0 → CERW.IsSubgradient Ψ (CERW.toSpace x) (ξ x)) → ξ 0 = 0 →
      ∀ n : ℕ, 2 ≤ n → ∀ q : EuclideanSpace ℝ (Fin d), ‖q‖ ≤ CERW.normMax Ψ →
      ∀ x : ℕ → Site d, x 0 = 0 → (∀ j, x (j + 1) - x j ∈ CERW.unitSteps d) →
        let Λ : ℝ := CERW.normMax Ψ
        let Z : ℕ → EuclideanSpace ℝ (Fin d) := fun t => CERW.toSpace (x t) + ε •
          ∑ j ∈ Finset.range t, if x j ∉ CERW.departureRange x j then ξ (x j) else 0
        (∀ s t : ℕ, s < t → t ≤ n →
          ‖Z t - Z s‖ ≤ C₁ * Real.sqrt ((t - s : ℝ) * Real.log n)) →
        (∀ k : ℕ, k ≤ n →
          |CERW.dynkinMart (CERW.driftStepProb d ε ξ)
              (fun z => max (inner ℝ q (CERW.toSpace z) - k * Λ) 0) x n|
            ≤ C₁ * (Real.sqrt (Real.log n * ∑ z ∈ (CERW.departureRange x n).filter
                  (fun z => k * Λ - Λ < inner ℝ q (CERW.toSpace z)),
                  (CERW.localTime x n z : ℝ)) + Real.log n)) →
        ∀ j₀ : ℕ, j₀ ≤ n →
        (∀ j ≤ n, inner ℝ q (CERW.toSpace (x j)) ≤ inner ℝ q (CERW.toSpace (x j₀))) →
        let T : ℝ := inner ℝ q (CERW.toSpace (x j₀))
        ∀ h : ℝ, 0 < h →
        (∀ j ≤ n, T - h < inner ℝ q (CERW.toSpace (x j)) → α ≤ inner ℝ q (ξ (x j))) →
        ∀ b : ℝ, 0 ≤ b →
          T ≤ max (T - h) b + C * (1 + ((((CERW.departureRange x n).filter
              (fun z => b ≤ inner ℝ q (CERW.toSpace z))).sup (CERW.localTime x n) : ℕ) : ℝ))
            * Real.log n
--
\end{verbatim}
\paragraph{Transcription.}
For d>=2, a norm, positive drift, coordinate ellipticity and p>0, the finite radial-vector/level family and vector event are produced with one event constant before all fields and laws. The separate deterministic implication quantifies every alpha>0 and chooses its bound before xi,n,q and the legal path; |q|<=normMax Psi, maximality, the vector bound, every source grid-level ramp bound and the drift-alpha condition are explicit. It concludes the original T bound for all b>=0. No simultaneous arbitrary-real-q probability event is claimed. The legal carrier is explicit; arbitrary IsDriftCERW paths are legal almost surely. Complete contract: ledger/nl/lem-outer-crossing.tex.
